\documentclass[11pt,a4paper]{article}
\usepackage[margin=1in]{geometry}
\usepackage{amsmath}
\usepackage{amssymb}
\usepackage{booktabs}
\usepackage{graphicx}
\usepackage{listings}
\usepackage{xcolor}
\usepackage{hyperref}
\usepackage{microtype}

\lstset{
  basicstyle=\ttfamily\small,
  breaklines=true,
  frame=single,
  backgroundcolor=\color{gray!10},
}

\title{State Pack: A Content-Addressed Protocol for\\Stateless Transformer Inference}

\author{%
  M.~Sadiq \\
  State Pack Project \\
  \texttt{https://github.com/mauludsadiq/State-Pack} \\
}

\begin{document}
\maketitle

\begin{abstract}
We present \textbf{State Pack}, a content-addressed protocol for efficient
transformer inference in multi-step agent workloads. State Pack treats KV cache
state as a portable, verifiable artifact rather than ephemeral server-side session
data. The core abstraction is a pure function:
$(\texttt{state\_hash},\ \texttt{delta}) \rightarrow (\texttt{new\_state\_hash},\ \texttt{output})$,
where \texttt{state\_hash} is the SHA-256 address of a serialized
\texttt{past\_key\_values} tensor. This design eliminates redundant prefill
computation across agent steps, producing 90--95\% token reduction across
GPT-2 (124M), Qwen2.5-3B, and Mistral-7B-Instruct on 20-step agent loops.
A split Rust/Python implementation achieves 9.1ms protocol latency alongside
47ms model inference on CPU. All results are fully reproducible:
\texttt{pip install state-pack}.
\end{abstract}

%-----------------------------------------------------------------------
\section{Introduction}
\label{sec:intro}

Autoregressive transformer models process tokens sequentially. At inference
time, the attention mechanism requires access to all prior key-value pairs ---
the KV cache. In agent loops, this cache grows with every step: a 40-step
loop with a 500-token base prompt and 30-token deltas requires processing
$[500, 530, 560, \ldots, 1670]$ tokens at steps 1 through 40. Total naive
cost: $\sim$43,000 tokens. Total \emph{necessary} computation: $\sim$1,700
tokens (base once, plus deltas).

KV cache state computed for a given input prefix is deterministic and
reusable. If the same base prompt is processed by the same model, the
resulting \texttt{past\_key\_values} tensor is identical. This makes KV
state \emph{content-addressable}: the SHA-256 hash of the input text
uniquely identifies the KV state it produces.

\subsection{Contributions}

\begin{enumerate}
  \item A formal stateless inference protocol:
        $(\texttt{state\_hash}, \delta) \rightarrow (\texttt{new\_state\_hash}, y)$
  \item Content-addressed KV cache serialization with float16 quantization
        (50\% size reduction, zero top-1 accuracy loss, verified experimentally)
  \item A \textsc{Compact} primitive preventing unbounded delta chain growth
  \item A \textsc{Gc} command for bounded storage in production deployments
  \item A split Rust/Python implementation: 9.1ms protocol latency, 12$\times$
        faster than the Python baseline
  \item Empirical evaluation across three model architectures showing 90--95\%
        token reduction
\end{enumerate}

%-----------------------------------------------------------------------
\section{Related Work}
\label{sec:related}

\paragraph{Prefix caching in inference engines.}
vLLM's PagedAttention~\cite{kwon2023} manages KV cache memory efficiently
but does not expose portable, addressable cache state across requests.
SGLang's RadixAttention~\cite{zheng2023} reuses prefixes within a single
server instance. LMCache~\cite{lmcache2024} supports KV offloading to
CPU/disk/S3 within vLLM but requires engine-level integration. State Pack
differs by treating KV state as a portable protocol artifact independent of
any specific inference engine.

\paragraph{Prompt caching in managed APIs.}
OpenAI, Anthropic, and Google offer prompt caching with 50--75\% discounts
on cached input tokens. These are server-side optimizations invisible to the
client. State Pack operates at the client layer, producing 90\%+ reduction
by avoiding transmission of redundant history entirely.

\paragraph{Content-addressed storage.}
IPFS~\cite{benet2014} and Nix~\cite{dolstra2006} demonstrate
content-addressable storage as infrastructure primitives. State Pack applies
this pattern to ML inference state, enabling deduplication, auditability, and
reproducibility not present in session-based inference systems.

%-----------------------------------------------------------------------
\section{Protocol Design}
\label{sec:protocol}

\subsection{State Representation}

A \emph{state packet} consists of:
\begin{itemize}
  \item \texttt{state\_hash}: $\text{SHA-256}(\texttt{base\_text})$
  \item \texttt{past\_key\_values}: serialized KV cache tensor (float16, \texttt{.pt})
  \item \texttt{manifest}: JSON metadata (\texttt{model}, \texttt{base\_tokens},
        \texttt{blob\_sha256}, \texttt{packet\_id})
  \item \texttt{packet\_id}: $\text{SHA-256}(\texttt{version} \mid \texttt{model}
        \mid \texttt{base\_sha256} \mid \texttt{base\_bytes} \mid
        \texttt{blob\_sha256} \mid \texttt{blob\_bytes})$
\end{itemize}

\subsection{Core Operations}

\begin{description}
  \item[\textsc{Create}$(t)$] Run base text $t$ through model, serialize
    \texttt{past\_key\_values}, store blob. Returns $h = \text{SHA-256}(t)$.
    \emph{Idempotent}: same text always produces same hash.

  \item[\textsc{Infer}$(h, \delta)$] Load KV state for $h$. Run $\delta$
    against cached state. Returns
    $h' = \text{SHA-256}(h \mid \text{SHA-256}(\delta))$ and output token.
    Emits tamper-evident receipt.

  \item[\textsc{Compact}$(h, [\delta_1,\ldots,\delta_n])$] Fold accumulated
    delta list into fresh base state. Returns $h_\text{fresh}$.
    Prevents unbounded growth of the delta chain.

  \item[\textsc{Gc}$(\texttt{store}, \tau)$] Remove manifest+blob pairs
    older than $\tau$ days or exceeding a maximum count.
    Keeps storage bounded in production deployments.
\end{description}

\subsection{Hash Chain}

The client maintains a hash chain:
\[
  h_0 \xrightarrow{\delta_1} h_1 \xrightarrow{\delta_2} h_2
  \xrightarrow{\phantom{xx}} \cdots
  \xrightarrow{\textsc{Compact}} h_\text{fresh}
\]
Each transition is deterministic given the same model weights. The chain
constitutes a cryptographic audit trail of inference history.

\subsection{Compaction Policy}

Without compaction, accumulated delta tokens eventually dominate base tokens,
reducing marginal savings. We define two policies:

\begin{itemize}
  \item \textbf{ThresholdPolicy}: compact when
        $\sum \texttt{delta\_tokens} > \texttt{base\_tokens} \times r$
  \item \textbf{SavingsPolicy}: compact when per-step savings fall below
        a minimum percentage
\end{itemize}

Empirically, \texttt{ThresholdPolicy}$(r{=}1.0, \texttt{max\_steps}{=}20)$
maintains $>$90\% savings over 20-step loops with 2 compactions on average.

%-----------------------------------------------------------------------
\section{Implementation}
\label{sec:impl}

\subsection{KV Cache Serialization}

HuggingFace \texttt{past\_key\_values} is a tuple of (key, value) tensor
pairs, one per layer. We serialize via \texttt{torch.save} with float16
quantization. Experimental verification (Section~\ref{sec:eval}) confirms
zero top-1 accuracy impact across all tested architectures.

\texttt{DynamicCache} compatibility (transformers $\geq$ 4.36) is handled
via \texttt{DynamicCache.from\_legacy\_cache()} and
\texttt{cache.to\_legacy\_cache()}.

\subsection{Split Architecture}

The protocol layer is implemented in Rust (Axum framework):
\begin{itemize}
  \item SHA-256 content addressing and manifest I/O
  \item Tamper-evident receipt generation
  \item HTTP routing
\end{itemize}

The inference layer is implemented in Python:
\begin{itemize}
  \item Model loading (HuggingFace \texttt{AutoModelForCausalLM})
  \item KV cache forward pass
  \item Blob serialization/deserialization
\end{itemize}

This split eliminates Python GIL contention on the protocol hot path.
Protocol latency: 9.1ms (Rust) vs.\ 107ms (Python baseline), 12$\times$
improvement.

\subsection{Multi-Agent Deduplication}

Content addressing naturally deduplicates shared base states. When $N$ agents
share the same system prompt, \textsc{Create} is called once and $N-1$
subsequent calls are cache hits. Empirically: first agent 0.951s, subsequent
agents 0.003s (316$\times$ faster). At 1,000 concurrent agents sharing a base
prompt, the cost reduction compound: base tokens paid once, delta tokens paid
per agent per step.

%-----------------------------------------------------------------------
\section{Evaluation}
\label{sec:eval}

\subsection{Token Reduction}

Table~\ref{tab:tokens} reports token reduction across model families and
sizes. Savings are consistent at 90--95\%, confirming the reduction is
structural (protocol-driven) rather than model-dependent.

\begin{table}[h]
\centering
\caption{Token reduction across model families (20-step agent loops).}
\label{tab:tokens}
\begin{tabular}{lrrrr}
\toprule
Model & Params & Naive Tokens & SP Tokens & Reduction \\
\midrule
GPT-2            & 124M & 18,660 & 875   & 95.3\% \\
Qwen2.5-3B       & 3B   & 5,412  & 495   & 90.9\% \\
Mistral-7B       & 7B   & 6,662  & 605   & 90.9\% \\
OpenAI API       & ---  & 17,929 & 1,320 & 92.6\% \\
\bottomrule
\end{tabular}
\end{table}

\subsection{Quantization Accuracy}

Float16 vs.\ float32 KV cache comparison (GPT-2):
\begin{itemize}
  \item Max logit difference: 0.066 (float16), 0.307 (bfloat16)
  \item Top-1 token match: \textbf{True} (both precisions)
  \item Blob size reduction: 50\%
\end{itemize}

\subsection{Protocol Latency}

\begin{table}[h]
\centering
\caption{Protocol layer latency: Rust vs.\ Python.}
\label{tab:latency}
\begin{tabular}{lrrr}
\toprule
Implementation & Avg & Min & Max \\
\midrule
Python (baseline) & 107ms & 24ms & 190ms \\
Rust (v0.3)       & 9.1ms & 2.8ms & 19.7ms \\
\midrule
Improvement & $12\times$ & --- & --- \\
\bottomrule
\end{tabular}
\end{table}

\subsection{Compaction Effectiveness}

\texttt{ThresholdPolicy}$(r{=}0.5, \texttt{max\_steps}{=}10)$ over 20 steps:
compactions triggered at steps 4 and 14, final savings 90.8\%,
average step latency 42ms.

\subsection{Limitations}

\begin{enumerate}
  \item \textbf{CPU-bound speedup.} Wall-clock speedup is 1.4$\times$ for 7B
        models on CPU. GPU inference is expected to yield 3--4$\times$
        improvement based on GPT-2 results; this remains future work.
  \item \textbf{KV cache portability.} Cache blobs are architecture-specific.
        A Mistral-7B blob cannot be loaded into a Qwen2.5 model.
        Cross-device transfer within the same architecture is not yet
        implemented.
  \item \textbf{OpenAI integration mechanism.} The OpenAI benchmark uses
        structured context discipline (delta-only messages) rather than
        KV cache transfer, which is unavailable via the OpenAI API.
        Native KV transfer would further improve savings.
  \item \textbf{Non-determinism.} Receipts hash inputs only. With
        temperature $> 0$, the same receipt may correspond to different
        outputs. Auditable deployments should use greedy decoding.
\end{enumerate}

%-----------------------------------------------------------------------
\section{Discussion}
\label{sec:discussion}

\paragraph{When State Pack helps most.}
Long agent loops ($>$10 steps) with stable base prompts; multi-agent
workloads with shared system prompts; regulated environments requiring
auditable inference; cost-sensitive deployments on frontier models.

\paragraph{When State Pack helps less.}
Single-turn queries with no prefix reuse; highly variable base prompts;
inference providers with aggressive native prompt caching that already
eliminates redundant prefill cost.

\paragraph{Relationship to existing systems.}
State Pack is complementary to vLLM and SGLang, not competitive. A complete
production deployment could use State Pack for client-side protocol and KV
addressability, with vLLM serving actual inference requests.

%-----------------------------------------------------------------------
\section{Conclusion}
\label{sec:conclusion}

State Pack demonstrates that treating transformer KV cache as a
content-addressed, portable protocol artifact produces consistent 90--95\%
token reduction across model families from 124M to 7B parameters. The
stateless pure-function design enables horizontal scalability, natural
multi-agent deduplication, and cryptographic auditability without
engine-level modifications. A Rust protocol implementation achieves 9.1ms
overhead, making the protocol layer negligible relative to model inference
time. All results are reproducible from a single command.

\begin{thebibliography}{9}

\bibitem{kwon2023}
W.~Kwon et al.,
\emph{Efficient Memory Management for Large Language Model Serving with
PagedAttention},
SOSP 2023.

\bibitem{zheng2023}
L.~Zheng et al.,
\emph{SGLang: Efficient Execution of Structured Language Model Programs},
2023.

\bibitem{lmcache2024}
LMCache Team,
\emph{LMCache: KV Cache Offloading for vLLM},
2024.

\bibitem{benet2014}
J.~Benet,
\emph{IPFS --- Content Addressed, Versioned, P2P File System},
2014.

\bibitem{dolstra2006}
E.~Dolstra,
\emph{The Purely Functional Software Deployment Model},
PhD thesis, Utrecht University, 2006.

\end{thebibliography}

\end{document}
